% --- Archivo: 01_clasificacion_convergencia.tex ---

La optimización matemática continua estudia el problema de minimizar una función sobre un dominio delimitado por restricciones. Representamos este modelo mediante la estructura general:
\begin{equation}
    \min f(x) \quad \text{sujeto a} \quad x \in D = \{x \in \Rn \mid h(x) = 0, \, g(x) \le 0\},
    \label{eq:v2_general_opt_problem}
\end{equation}
donde $\Rn$ es el espacio base, la función objetivo es $f : \Rn \to \R$, y las restricciones de igualdad y desigualdad vienen dadas por los operadores $h : \Rn \to \R^l$ y $g : \Rn \to \R^m$, respectivamente.

La resolución numérica de \eqref{eq:v2_general_opt_problem} requiere evaluar de forma sistemática los valores de la función objetivo, de las restricciones y, de forma habitual, de sus respectivas derivadas. Para estructurar el estudio teórico de estos algoritmos, presentaremos primero una clasificación sistemática de los métodos numéricos de optimización.

% =========================================================
\section{Clasificación de los métodos y nociones de convergencia}
\label{sec:v2_clasificacion_metodos}
% =========================================================

Los métodos de optimización se clasifican atendiendo a tres criterios fundamentales: la dependencia de la información obtenida en iteraciones previas (métodos pasivos o secuenciales), la cantidad de información histórica acumulada (esquemas con o sin memoria) y el orden máximo de las derivadas utilizadas.

\subsection{Estrategia de muestreo: Métodos pasivos vs. secuenciales}

Atendiendo al papel que desempeña la información obtenida durante el proceso de cálculo, distinguimos dos clases de métodos:
\begin{itemize}
    \item \textbf{Método pasivo:} La secuencia de puntos de evaluación se elige y predetermina completamente a priori, de forma que la elección de cada candidato no se beneficia de la información obtenida en las iteraciones previas (por ejemplo, la búsqueda en rejilla o \textit{grid search}).
    \item \textbf{Método secuencial:} Cada nuevo punto de la secuencia se determina utilizando los valores de la función objetivo, las restricciones o sus derivadas calculados en pasos anteriores. Esta clase resulta considerablemente más eficiente y abarca a la gran mayoría de los algoritmos computacionales prácticos.
\end{itemize}

\subsection{Estructura del esquema iterativo: Métodos sin memoria vs. con memoria}

Un método secuencial genera una secuencia de puntos en $\Rn$, denotada por $\{x^1, x^2, \dots, x^k, \dots\}$, a los que se denomina \deftech{aproximaciones} o \deftech{iterados} del método. Esta secuencia (llamada también \textbf{trayectoria}) se describe mediante un esquema iterativo.

Si denotamos por $\Psi_k$ al operador de transición definido sobre conjuntos $U_k \subset \Rn$, la generación de un nuevo iterado se expresa como:
\begin{equation}
    x^{k+1} = \Psi_k(x^k), \quad k = 0, 1, \dots
    \label{eq:v2_iterative_scheme_base}
\end{equation}
Para que el esquema esté bien definido a partir de un punto inicial $x^0 \in U_0$, los sucesivos iterados deben permanecer dentro de los dominios de definición de los operadores correspondientes. Esta condición de compatibilidad se expresa como:
\[
    \Psi_{k-1}(\Psi_{k-2}(\dots \Psi_0(x^0)\dots)) \in U_k \quad \forall k = 1, 2, \dots
\]
Dependiendo de la cantidad de información histórica que requiera el operador $\Psi_k$, clasificamos los esquemas iterativos en:
\begin{itemize}
    \item \textbf{Esquema sin memoria:} El operador $\Psi_k$ depende únicamente de la aproximación actual $x^k$.
    \item \textbf{Esquema con memoria:} El operador utiliza información del iterado actual $x^k$ y de puntos anteriores de la trayectoria ($x^{k-1}, x^{k-2}, \dots$).
\end{itemize}

\subsection{Orden del método}

El tercer criterio clasifica los algoritmos según la información diferencial que emplean.

\begin{definition}[Orden de un método]\label{def:v2_method_order}
    El \deftech{orden de un método} es el orden máximo de las derivadas de la función objetivo y de las restricciones utilizadas en cada iteración del algoritmo.
    \begin{enuthm}
        \item \textbf{Métodos de orden cero (o directos):} Aquellos que solo evalúan funcionalmente $f$, $g$ o $h$, sin utilizar derivadas.
        \item \textbf{Métodos de primer orden:} Aquellos que basan sus decisiones en la evaluación de las primeras derivadas (vectores gradientes, $\nabla f$).
        \item \textbf{Métodos de segundo orden:} Aquellos que requieren el cálculo de segundas derivadas (matrices Hessianas, $\nabla^2 f$).
    \end{enuthm}
\end{definition}

% ---------------------------------------------------------
\subsection{Tipos de convergencia}
% ---------------------------------------------------------

El comportamiento asintótico de la trayectoria $\{x^k\}$ define las propiedades de convergencia del algoritmo. Distinguimos entre \deftech{convergencia finita} (cuando la trayectoria alcanza una solución exacta en un número finito de pasos, es decir, $x^k \in S$ para algún $k < \infty$) y \deftech{convergencia asintótica}.

En problemas no lineales generales, el análisis se centra rigurosamente en la convergencia asintótica. Dependiendo de la métrica y del espacio considerados, distinguimos cuatro tipos de convergencia:

\begin{enuthm}
    \item \textbf{Convergencia en relación a las variables (o puntual):} La sucesión de iterados converge a una solución del problema:
    \[
        x^k \to \bar{x} \quad (k \to \infty), \quad \text{con } \bar{x} \in S.
    \]
    Esta propiedad suele requerir hipótesis restrictivas (como convexidad estricta o unicidad de solución) para poder ser demostrada de forma directa.

    \item \textbf{Convergencia al conjunto de soluciones $S$:} La distancia de la trayectoria al conjunto de soluciones tiende a cero:
    \[
        \dist(x^k, S) \to 0 \quad (k \to \infty).
    \]
    Si el conjunto $S$ es compacto, esto asegura que el conjunto de puntos de acumulación de $\{x^k\}$ es no vacío y está contenido en $S$. Si $S$ es ilimitado, es posible que la distancia tienda a cero sin que existan puntos de acumulación.

    \item \textbf{Convergencia en relación a la función objetivo (o al valor óptimo):} La sucesión de valores de la función converge al valor mínimo $\bar{v}$, garantizando simultáneamente la aproximación al conjunto factible $D$:
    \[
        \dist(x^k, D) \to 0 \quad \text{y} \quad f(x^k) \to \bar{v} \quad (k \to \infty).
    \]
    Este tipo de convergencia no implica necesariamente que la trayectoria de las variables $\{x^k\}$ sea convergente en el espacio $\Rn$.
    
    \item \textbf{Convergencia en relación al gradiente:} En optimización irrestricta con funciones no convexas, es muy difícil garantizar la convergencia hacia un mínimo global. Por ello, el análisis se suele centrar en demostrar que la norma del gradiente converge a cero:
    \[
        \norm{\nabla f(x^k)} \to 0 \quad (k \to \infty).
    \]
    Bajo esta condición, cualquier punto de acumulación $\bar{x}$ de la trayectoria es, por lo menos, un punto estacionario del problema ($\nabla f(\bar{x}) = 0$).
\end{enuthm}

Distinguimos también entre \textbf{convergencia global} y \textbf{convergencia local}. Un método tiene \textit{convergencia global} si la trayectoria converge a una solución para cualquier punto inicial $x^0$. Por el contrario, un método \textit{local} requiere que $x^0$ se sitúe suficientemente cerca de la solución para garantizar su convergencia.

Frecuentemente, la resolución de un problema complejo se basa en reducirlo a una secuencia de problemas más simples. Por ejemplo, el método de Newton aproxima un sistema no lineal mediante sistemas lineales. De igual modo, muchos algoritmos multidimensionales se apoyan internamente en la resolución iterativa de problemas unidimensionales a lo largo de direcciones específicas (búsqueda lineal).

\begin{notabox}[El problema de la Optimización Global]
    Para problemas fuertemente no convexos, hallar el mínimo global es, en general, analíticamente imposible. El enfoque ingenuo de evaluar una red muy fina sobre el conjunto factible sufre la llamada \textit{maldición de la dimensionalidad}.
    
    Un enfoque heurístico más eficiente es la estrategia de \textbf{múltiple partida (multi-start)}. Esta estrategia superpone una red discreta gruesa sobre el conjunto factible y ejecuta un método de optimización local rápido desde cada nodo. La aproximación a la solución global se obtiene comparando los valores de los óptimos locales alcanzados por cada trayectoria.
\end{notabox}